% !TEX root = ../main.tex
\chapter*{Resum}
\addcontentsline{toc}{chapter}{Resum}
Els videojocs d'estratègia amb informació imperfecta, accions simultànies i dinàmica estocàstica constitueixen un banc de proves exigent per a la presa de decisions en Intel·ligència Artificial. Aquest Treball de Final de Màster formula, desenvolupa i avalua cinc paradigmes d'agents en el format individual (\texttt{gen9randombattle}) de Pokémon competitiu, utilitzant el simulador Pokémon Showdown i el marc \textit{poke-env}.

El problema es formalitza com un joc estocàstic d'observació parcial (POSG) de dos jugadors i suma zero amb accions simultànies, reduït a POMDP en PPO i XGBoost en operar sobre observacions parcials $o_i \in \Omega_i$. En un espai combinatòriament ingent amb equips i moviments ocults, l'objectiu no és només maximitzar la taxa de victòries sota pressupostos explícits, sinó analitzar, mitjançant telemetria per torn, la gestió del risc, la predicció de canvis, el ritme i l'ús estratègic de la Teracristal·lització.

S'implementen cinc paradigmes homogenis: (1)~una seqüència d'ablació de catorze heurístiques (\texttt{v1}--\texttt{v14}) amb coneixement de domini incremental (càlcul de dany, efectivitat de tipus, \emph{hazards}, preparació, estats alterats i Teracristal·lització); (2)~cerca adversària d'un torn (\emph{minimax} \texttt{v15}--\texttt{v17}) amb funcions d'avaluació a les fulles; (3)~cerca en arbre de Monte Carlo superficial (\emph{shallow} MCTS \texttt{v18}--\texttt{v20}, 5 torns) comparant UCB1 i PUCT amb priors heurístics; (4)~aprenentatge per imitació (\texttt{v21}--\texttt{v22}) amb XGBoost sobre partides humanes d'alt nivell (Elo~$1800+$) per classificar atacs o canvis; i (5)~aprenentatge per reforç profund (\emph{MaskablePPO}) amb accions emmascarades i clonatge conductual. L'avaluació empra una infraestructura paral·lela amb telemetria per torn i rànquings Bradley--Terry.

L'avaluació empírica sobre 28 agents (6,4 milions de batalles) i proves de DRL situa l'heurística \texttt{v12} com a líder global, superant el referent competitiu \textit{Abyssal} amb un 59,9\% de victòries. \texttt{v14}, més exhaustiva en regles i utilitzada com a fulla, prior i agent al \emph{ladder} públic, no millora el rendiment de \texttt{v12}. Cap variant de cerca, imitació o DRL supera \texttt{v12} sota els pressupostos avaluats. La telemetria conclou que la incertesa de l'estat ocult rival limita la cerca superficial, mentre que els mètodes apresos requereixen integrar coneixement de domini expert per assolir nivell competitiu.

\vspace{0.3cm}
\noindent\textbf{Paraules clau:} Intel·ligència Artificial en jocs, presa de decisions sota incertesa, POSG, POMDP, agents heurístics, cerca adversària (Minimax), cerca en arbre de Monte Carlo (MCTS), aprenentatge per imitació (XGBoost), aprenentatge per reforç profund (MaskablePPO), telemetria, perills d'entrada (\emph{hazards}), model Bradley--Terry.

\clearpage

\chapter*{Resumen}
\addcontentsline{toc}{chapter}{Resumen}
Los videojuegos de estrategia con información imperfecta, acciones simultáneas y dinámica estocástica plantean un reto complejo para la toma de decisiones en Inteligencia Artificial. Este Trabajo de Fin de Máster formula, desarrolla y evalúa cinco paradigmas de agentes en el formato individual (\texttt{gen9randombattle}) del videojuego competitivo Pokémon, utilizando Pokémon Showdown y el marco \textit{poke-env}.

El problema se formaliza como un juego estocástico de observación parcial (POSG) de dos jugadores y suma cero con acciones simultáneas, reducido a POMDP en PPO y XGBoost al operar sobre observaciones locales $o_i \in \Omega_i$. En un espacio combinatoriamente inmenso con equipos y movimientos ocultos, el objetivo no es solo maximizar victorias bajo presupuestos explícitos, sino analizar mediante telemetría por turno la gestión del riesgo, la predicción de cambios, el ritmo y el uso táctico de la Teracristalización.

Se implementan cinco paradigmas homogéneos: (1)~una serie de ablación de catorce heurísticas (\texttt{v1}--\texttt{v14}) con conocimiento incremental (cálculo de daño, efectividad de tipos, \emph{hazards}, potenciación, estados y Teracristalización); (2)~búsqueda adversaria de un turno (\emph{minimax} \texttt{v15}--\texttt{v17}) con evaluación en hojas; (3)~búsqueda en árbol de Monte Carlo superficial (\emph{shallow} MCTS \texttt{v18}--\texttt{v20}, 5 turnos) contrastando UCB1 y PUCT con priors heurísticos; (4)~aprendizaje por imitación (\texttt{v21}--\texttt{v22}) con XGBoost sobre partidas humanas de élite (Elo~$1800+$) para clasificar ataques o cambios; y (5)~aprendizaje por refuerzo profundo (\emph{MaskablePPO}) con acciones enmascaradas y clonación conductual. La evaluación utiliza una infraestructura paralela con telemetría por turno y clasificaciones Bradley--Terry.

La evaluación empírica sobre 28 agentes (6,4 millones de batallas) y pruebas de DRL posiciona a la heurística \texttt{v12} como líder global, superando al referente competitivo \textit{Abyssal} con un 59,9\% de victorias. \texttt{v14}, más exhaustiva en reglas y empleada como hoja, prior y agente en el \emph{ladder} público, no supera a \texttt{v12}. Ninguna variante de búsqueda, imitación o DRL supera a \texttt{v12} bajo los presupuestos evaluados. La telemetría revela que la incertidumbre del estado oculto rival limita la búsqueda superficial, y que los modelos aprendidos requieren incorporar conocimiento experto de dominio para ser competitivos.

\vspace{0.3cm}
\noindent\textbf{Palabras clave:} Inteligencia Artificial en juegos, toma de decisiones bajo incertidumbre, POSG, POMDP, agentes heurísticos, búsqueda adversaria (Minimax), búsqueda en árbol de Monte Carlo (MCTS), aprendizaje por imitación (XGBoost), aprendizaje por refuerzo profundo (MaskablePPO), telemetría, peligros de entrada (\emph{hazards}), modelo Bradley--Terry.

\clearpage

\chapter*{Abstract}
\addcontentsline{toc}{chapter}{Abstract}
Complex strategy video games with imperfect information, simultaneous actions, and stochastic dynamics represent a demanding testbed for decision-making in Artificial Intelligence. This Master's Thesis formulates, implements, and benchmarks five agent paradigms for the singles format (\texttt{gen9randombattle}) of competitive Pokémon, leveraging Pokémon Showdown and the \textit{poke-env} framework.

The environment is formalized as a two-player, zero-sum Partially Observable Stochastic Game (POSG) with simultaneous actions, reduced to a POMDP in PPO and XGBoost via partial observations $o_i \in \Omega_i$. Across a combinatorially vast state space with hidden teams and movepools, the objective extends beyond maximizing win rates under explicit budgets: fine-grained per-turn telemetry is used to analyze risk management, switch prediction, tempo control, and tactical Terastallization.

Five standardized paradigms are developed: (1)~a 14-agent ablation sequence (\texttt{v1}--\texttt{v14}) incrementally incorporating domain expertise (damage formulas, type charts, hazards, setup moves, status conditions, and Terastallization); (2)~one-turn adversarial search (\emph{minimax} \texttt{v15}--\texttt{v17}) with leaf evaluations; (3)~shallow Monte Carlo Tree Search (\emph{shallow} MCTS \texttt{v18}--\texttt{v20}, 5-turn rollouts) comparing UCB1 with heuristic-prior-guided PUCT; (4)~imitation learning (\texttt{v21}--\texttt{v22}) using XGBoost trained on high-level human replays (Elo~$1800+$) to predict attack vs. switch decisions; and (5)~Deep Reinforcement Learning (\emph{MaskablePPO}) with action masking and Behavioral Cloning initialization. Evaluation relies on a parallel, resilient pipeline with per-turn telemetry and Bradley--Terry rankings.

Empirical benchmarking across 28 agents (6.4 million battles) and dedicated DRL experiments identifies heuristic \texttt{v12} as the top performer, surpassing the competitive baseline \textit{Abyssal} with a 59.9\% win rate. Meanwhile, \texttt{v14}—the most rule-complete agent, used as leaf evaluator, search prior, and live ladder agent—does not improve upon \texttt{v12}. Under the evaluated computational budgets, no search, imitation, or DRL variant beats \texttt{v12}. Telemetry demonstrates that hidden state uncertainty severely constrains shallow search, while learned models strictly require expert domain knowledge to achieve competitive play.

\vspace{0.3cm}
\noindent\textbf{Keywords:} Game Artificial Intelligence, decision-making under uncertainty, POSG, POMDP, heuristic agents, adversarial search (Minimax), Monte Carlo Tree Search (MCTS), imitation learning (XGBoost), Deep Reinforcement Learning (MaskablePPO), telemetry, entry hazards, Bradley--Terry model.
